\NumberedItem{7}
\HistoricTitle{\textsc{Gauss} an \textsc{Laplace}.}
\begin{flushright}Göttingen le 5 nov. 1812\end{flushright}
\noindent\hspace*{.27\linewidth}Monsieur

\indent Je vous dois encore mille remercimens pour le precieux cadeau, que vous m’avés fait l’honneur de me faire de votre \textit{Théorie des Probabilités}[\textsuperscript{**})]: M. \textsc{Olbers} m’a remis le second volume et Mr. \textsc{de Lindenau} le premier. Il est superflu de dire, combien la lecture de cet ouvrage m’a été et est encore intéressante et instructive. Vous aurés reçu une copie d’un mémoire sur une fonction transcendante très-générale, que j’ai eu l’honneur de vous envoïer. Ce mémoire ne contient que la première partie de mes recherches sur cette matière, je n’ai pas encore eu le loisir d’achever le second mémoire, qui en contiendra la suite. En effet outre les occupations habituelles de ma place, les calculs sur les perturbations m’ont couté depuis deux ans beaucoup de tems. Souvent je serais tenté de le regretter comme pouvant [être] peut être mieux
\par\medskip\noindent\rule{.24\linewidth}{.4pt}\par\smallskip
{\small
\noindent\hspace*{\parindent}[\textsuperscript{**}) P. S. \textsc{Laplace}, \textit{Théorie analytique des Probabilités}, Paris 1812.]
\par}
\clearpage\SourceLeaf{386}\ArithmeticHead{379}{R}{GAUSS AN LAPLACE.}
\noindent emploïé qu’à des calculs machinals, si la considération de l’utilité d’un tel travail ne me consolait. Peut être il faudra encore une année entière pour achever les calculs de la perturbation produite par Jupiter seulement. Récemment pour prendre un peu de relâche après des calculs fastidieux, je me suis occupé du problème célèbre des attractions d’un spheroïde elliptique[\textsuperscript{*})]. C’est vous, Monsieur, qui il y a 30 ans en avés donné le premier la solution complète[\textsuperscript{**})], dont j’ai admiré tant de fois la subtilité. Je me flatte que la manière nouvelle, dont je traite cette question meritera l’attention des geometres. J’ai composé la dessus un mémoire qui sera lu bientôt à la Société Roïale et ensuite imprimé parmi ses »Commentationes«[\textsuperscript{***})]. J’ai l’honneur de vous offrir ici un Extrait de ce qui est essentiel au problème cité, et je vous prie de le presenter à l’Institut, duquel plusieurs membres ont bien merité du même problème. Vous verrés avec plaisir que deux pages m’ont suffi pour obtenir la solution complète.

\indent Continués Monsieur de m’honorer de votre bienveillance et agréés l’assurance de la plus profonde estime, avec laquelle j’ai l’honneur d’être votre très humble serviteur
\begin{flushright}\textsc{Charles-Frederic Gauss}.\end{flushright}
\par\medskip\noindent\rule{.24\linewidth}{.4pt}\par\smallskip
{\small
\noindent\hspace*{\parindent}[\textsuperscript{*}) Vergl. die Tagebuchaufzeichnungen Nr. 142 vom 26. September und Nr. 143 vom 15. Oktober 1812.]
\par\noindent\hspace*{\parindent}[\textsuperscript{**}) \textit{Théorie des attractions des sphéroïdes et de la figure des planètes}, Histoire de l’Académie des Sciences, Année 1782, Paris 1785, Mémoires de Math. et Phys., S. 113.]
\par\noindent\hspace*{\parindent}[\textsuperscript{***}) \textit{Theoria attractionis corporum sphaeroidicorum ellipticorum homogeneorum methodo novo tractata.} Soc. Reg. Scient. tradita XVIII. Mart. MDCCCXIII, Werke V, S. 1.]
\par}
